\documentclass[plain]{baustms}

%% The class already loads amsmath, amssymb, amsthm, txfonts, geometry,
%% enumerate and hyperref; no further packages are required.
%% The production class displays a submission banner and placeholder DOI by
%% default. Suppress both in this author-prepared submission manuscript.
\makeatletter
\renewcommand\submittedto{}
\def\@mainheadline{\gdef\@mymainhead{\mbox{}}\setcounter{page}{1}}
\makeatother

\newcommand{\C}{\mathbb C}
\newcommand{\N}{\mathbb N}
\newcommand{\Q}{\mathbb Q}
\newcommand{\D}{\mathbb D}
\newcommand{\Pp}{\mathbb P}
\newcommand{\E}{\mathbb E}
\newcommand{\ind}{\mathbf 1}
\newcommand{\Real}{\operatorname{Re}}

\theoremstyle{cupthm}
\newtheorem{theorem}{Theorem}[section]
\newtheorem{lemma}[theorem]{Lemma}
\newtheorem{corollary}[theorem]{Corollary}

\runningtitle{Cofinite Zeros of High Derivatives}
\authorheadline{E. Hou}

\title{Cofinite Zeros of High Derivatives}
\author{Eric Hou}
\address{Independent researcher, Los Gatos, California, USA%
\email{eric.x.hou@gmail.com}}

\begin{document}

\begin{abstract}
We construct a transcendental entire function $f$ for which every nonempty
open subset of the complex plane contains a zero of $f^{(n)}$ for all
sufficiently large~$n$. Earlier Gaussian proposals by Almeida and Chojecki use
expected-zero and Offord-type estimates. Here the coefficients are independent
and uniformly distributed on the closed unit disk. A saddle-point estimate,
one-coordinate anti-concentration and the mean value property of
$\log|f^{(n)}|$ give summable fixed-disk hole probabilities. Every sample
satisfies $|f(z)|\leq\sqrt2\exp(|z|^2)$. The selected function is therefore a
counterexample to Boas and Reddy's printed assertion that every transcendental
entire function of order at most two and finite type admits an arbitrarily
large fixed disk zero-free for infinitely many successive derivatives.
\end{abstract}

\classification[2020]{primary 30D20; secondary 30B20, 30C15, 60F15}
\keywords{entire function, successive derivatives, zero sets,
random analytic function, hole probability}

\maketitle

\section{Introduction}

Erd\H{o}s asked whether there is a nonzero entire function $f$ such that, for
every strictly increasing sequence $(n_k)_{k\geq0}$ of nonnegative integers,
the union of the zero sets of $f^{(n_k)}$ is dense in the complex plane
\cite[p.~72]{Erdos1982}. On the same page he wrote that the relevant functions
had been constructed ``more than ten years ago'', without identifying a
source. The question is now recorded as Problem~906 \cite{Bloom906}. If
polynomials are admitted it is immediate, so we treat the transcendental
form. Write $\N=\{0,1,2,\ldots\}$ and
$B(c,R)=\{z\in\C:|z-c|<R\}$.

For a nonempty open set $U\subset\C$ and an entire function $f$, define
\[
 A_U(f)=\{n\in\N:f^{(n)}\text{ has a zero in }U\}.
\]
The following elementary reformulation isolates the quantifier that the
problem requires.

\begin{lemma}\label{lem:cofinite}
For an entire function $f$, the following assertions are equivalent.
\begin{enumerate}[\rm(1)]
\item For every strictly increasing sequence $(n_k)_{k\geq0}$ in $\N$, the set
\[
 \bigcup_{k\geq0}\{z\in\C:f^{(n_k)}(z)=0\}
\]
is dense in $\C$.
\item For every nonempty open set $U\subset\C$, the set $A_U(f)$ is cofinite
in $\N$, that is, $\N\setminus A_U(f)$ is finite.
\end{enumerate}
\end{lemma}

\begin{proof}
Assume (2). Fix a strictly increasing sequence $(n_k)$ and a nonempty open
set $U$. Since $A_U(f)$ is cofinite there is an integer $N(U)$ such that
$n\geq N(U)$ implies $n\in A_U(f)$. The sequence $(n_k)$ is unbounded, so
$n_k\geq N(U)$ for some $k$, and then $f^{(n_k)}$ has a zero in $U$. The
displayed union therefore meets every nonempty open set and is dense.

Conversely, suppose that (2) fails. Then for some nonempty open set $U$ the
set $\N\setminus A_U(f)$ is infinite; enumerate it increasingly as $(n_k)$.
Each $f^{(n_k)}$ is zero-free in $U$, so the union of their zero sets misses
$U$ and is not dense.
\end{proof}

\begin{theorem}\label{thm:main}
There exists a transcendental entire function $f$ such that, for every
nonempty open set $U\subset\C$, there is an integer $N(U)$ for which
\[
 n\geq N(U)\quad\Longrightarrow\quad f^{(n)}\text{ has a zero in }U.
\]
The same function can be chosen to satisfy
\begin{equation}\label{eq:growth}
 |f(z)|\leq\sqrt2\exp(|z|^2),\qquad z\in\C.
\end{equation}
Consequently, the zero sets indexed by every strictly increasing sequence of
derivative orders have dense union.
\end{theorem}

The bound \eqref{eq:growth} implies that the function below has order at most
two and, if its order is two, type at most one. The Boas--Reddy consequence,
with its source quantifier made explicit, is recorded in
Section~\ref{sec:boas-reddy}.

MacLane \cite{MacLane1952} constructed an entire function whose successive
derivatives are dense in the space of entire functions, and Barth and
Schneider \cite{BarthSchneider1972} prescribed discrete zero sets along one
constructed sequence of derivative orders. Neither result yields the
cofinite condition in Lemma~\ref{lem:cofinite}. Boas and Reddy's expanded
article \cite{BoasReddyJMAA1973} contained a further positive Theorem~2, but
their erratum withdrew its construction and left that conclusion open
\cite{BoasReddy1975}.

On 25 April 2026, Almeida first posted a detailed proposed solution and later
uploaded a manuscript using the planar Gaussian entire function
\cite{AlmeidaForum2026,Almeida2026}. Later that day Chojecki posted a second
Gaussian-route note; on 1 June he clarified that it was intended as a full
solution and reiterated that Almeida had posted first
\cite{ChojeckiForum2026}. These posts precede the present manuscript. The
author learned of them after completing and circulating this proof.

Almeida's manuscript and the present work share the cofinite reformulation,
the Fock normalization $1/\sqrt{k!}$ and a final Borel--Cantelli argument.
Almeida uses Gaussian coefficients, the Edelman--Kostlan expected zero
measure, Laguerre asymptotics and an Offord-type estimate. Chojecki's forum
description likewise invokes the cofinite reformulation and the
Edelman--Kostlan--Offord Gaussian route. The contribution here is an
alternative bounded-coefficient proof with a different local hole-probability
mechanism and a deterministic finite-type growth bound, the latter yielding
the Boas--Reddy counterexample. Specifically, a saddle-selected coefficient,
one-coordinate anti-concentration and the mean value property for
$\log|f^{(n)}|$ give summable hole probabilities; Borel--Cantelli
\cite[Section~2.3]{Durrett} and a countable disk basis then produce one sample.
For background on random series, see Kahane \cite{Kahane}. As accessed 16
August 2026, the online record still displays \textsc{open} while recording
claimed solutions and warning that relevant literature may be missing
\cite{Bloom906,AlmeidaForum2026,ChojeckiForum2026}.

\section{A bounded random Fock series}\label{sec:fock}

Let $\overline\D$ be the closed unit disk and let $m_\D$ be normalised planar
Lebesgue measure on $\overline\D$. Put
\[
 \Omega=\overline\D^{\,\N},\qquad
 \mathcal F=\mathcal B(\overline\D)^{\otimes\N},\qquad
 \Pp=m_\D^{\otimes\N},
\]
and equip $\Omega$ with the product topology. A countable product of compact
metric spaces is compact and metrizable, and $\mathcal F$ is the Borel
$\sigma$-algebra of $\Omega$. For $\omega=(\omega_k)_{k\geq0}\in\Omega$ set
$\xi_k(\omega)=\omega_k$. The coordinate variables $\xi_0,\xi_1,\ldots$ are
then independent and uniformly distributed on $\overline\D$. Define
\begin{equation}\label{eq:fock}
 f(z)=\sum_{k=0}^{\infty}\xi_k\frac{z^k}{\sqrt{k!}}.
\end{equation}
The bound $|\xi_k|\leq1$ gives normal convergence on compact subsets of $\C$,
uniformly in $\omega$. Hence every sample defines an entire function, term by
term differentiation is legitimate, and for $n\geq0$,
\begin{equation}\label{eq:derivative}
 F_n(z):=f^{(n)}(z)
 =\sum_{m=0}^{\infty}\xi_{n+m}\frac{\sqrt{(n+m)!}}{m!}\,z^m.
\end{equation}
The partial sums of \eqref{eq:derivative} are continuous on
$\Omega\times\C$ and converge uniformly on $\Omega\times\{|z|\leq M\}$ for
every $M$: the absolute-value majorant has successive-term ratio
$M\sqrt{n+m+1}/(m+1)\to0$. Hence $(\omega,z)\mapsto F_n(\omega,z)$ is jointly continuous; in
particular $\omega\mapsto F_n(\omega,z)$ is $\mathcal F$-measurable for each
fixed $z$. The law $m_\D$ has no atoms, so $\Pp\{\xi_k=0\}=0$ for every $k$,
and therefore
\begin{equation}\label{eq:nonzero}
 \Omega_0=\{\omega\in\Omega:\xi_k(\omega)\neq0\text{ for every }k\geq0\}
\end{equation}
has $\Pp(\Omega_0)=1$. On $\Omega_0$ all Taylor coefficients in
\eqref{eq:fock} are nonzero, so $f$ is not a polynomial.

The bounded coefficients also give the growth bound. By Cauchy--Schwarz with
weights $2^{k/2}$,
\begin{align*}
 |f(z)|
 &\leq\sum_{k\geq0}\frac{|z|^k}{\sqrt{k!}}\\
 &\leq\left(\sum_{k\geq0}\frac{2^k|z|^{2k}}{k!}\right)^{1/2}
      \left(\sum_{k\geq0}2^{-k}\right)^{1/2}
  =\sqrt2\exp(|z|^2).
\end{align*}
Thus \eqref{eq:growth} holds for every $\omega\in\Omega$, and every nonzero
sample function has order at most two and finite type.

For $r\geq0$ write
\[
 a_{n,m}(r)=\frac{\sqrt{(n+m)!}}{m!}\,r^m,
 \qquad
 S_n(r)=\sum_{m=0}^{\infty}a_{n,m}(r),
\]
and set
\begin{equation}\label{eq:potential}
 p_n(z)=\tfrac12\log(n!)+|z|\sqrt n,
 \qquad
 L_n=2+\log(n+1).
\end{equation}
We use the extended-real convention $\log0=-\infty$. Every occurrence of
$L_n$ below is under hypotheses that force $n\geq1$.

The saddle-point estimate needs an upper bound for $\log(m!)$. Rather than
import a Stirling inequality we prove the following weaker bound, whose loss
relative to Stirling's estimate is one logarithm; only the property
$L_n=o(\sqrt n)$ is used, so the loss is harmless.

\begin{lemma}[Elementary factorial bound]\label{lem:factorial}
For every integer $m\geq1$,
\[
 \log(m!)\leq m\log m-m+\log(m+1)+1.
\]
\end{lemma}

\begin{proof}
The function $\log$ is nondecreasing on $[1,\infty)$, so
$\log j\leq\int_j^{j+1}\log t\,dt$ for every integer $j\geq1$. Summing over
$1\leq j\leq m$ and evaluating the integral,
\[
 \log(m!)=\sum_{j=1}^{m}\log j
 \leq\int_1^{m+1}\log t\,dt
 =\bigl[t\log t-t\bigr]_1^{m+1}
 =(m+1)\log(m+1)-m.
\]
Furthermore $(m+1)\log(m+1)=m\log(m+1)+\log(m+1)$, and since
$\log(1+u)\leq u$ for $u\geq0$,
\[
 m\log(m+1)=m\log m+m\log\Bigl(1+\frac1m\Bigr)\leq m\log m+1.
\]
Combining the two displays gives the assertion.
\end{proof}

\begin{lemma}[Saddle estimates]\label{lem:saddle}
Let $n\geq1$.
\begin{enumerate}[\rm(1)]
\item If $r\geq0$, $1\leq r\sqrt n$ and $r\leq\sqrt n$, then, for
$m=\lfloor r\sqrt n\rfloor$,
\[
 \log a_{n,m}(r)\geq\tfrac12\log(n!)+r\sqrt n-2-\log(n+1);
\]
by \eqref{eq:potential} the right-hand side equals $p_n(z)-L_n$ for every
$z$ with $|z|=r$.
\item For every $r\geq0$,
\[
 \log S_n(r)\leq\tfrac12\log(n!)+r\sqrt n+\frac{r+r^2}{2}.
\]
Consequently, if $|z|\leq R$ then
\begin{equation}\label{eq:upper}
 \log|F_n(z)|-p_n(z)\leq\frac{R+R^2}{2}
\end{equation}
for every $\omega\in\Omega$, with the inequality read trivially when
$F_n(z)=0$.
\end{enumerate}
\end{lemma}

\begin{proof}
Factor $a_{n,m}(r)=\sqrt{n!}\,b_{n,m}(r)$, where
\[
 b_{n,m}(r)=\frac{\sqrt{(n+1)(n+2)\cdots(n+m)}}{m!}\,r^m,
\]
which is legitimate because $(n+m)!=n!\,(n+1)(n+2)\cdots(n+m)$.

For (1) put $x=r\sqrt n$. Each of the $m$ factors under the square root is at
least $n$, so $b_{n,m}(r)\geq n^{m/2}r^m/m!=x^m/m!$. For $m=\lfloor x\rfloor$
the hypotheses give $1\leq m\leq x\leq n$, since $x\geq1$ and
$x=r\sqrt n\leq\sqrt n\sqrt n=n$. By Lemma~\ref{lem:factorial},
\[
 \log\frac{x^m}{m!}
 \geq m\log x-m\log m+m-\log(m+1)-1
 =m\log\frac xm+m-\log(m+1)-1.
\]
Here $m\log(x/m)\geq0$ because $m\leq x$; also $m>x-1$ and
$\log(m+1)\leq\log(n+1)$ because $m\leq n$. Hence
\[
 \log\frac{x^m}{m!}\geq(x-1)-\log(n+1)-1=x-2-\log(n+1),
\]
and adding $\tfrac12\log(n!)$ to both sides gives (1); the identification of
the right-hand side with $p_n(z)-L_n$ is the definition
\eqref{eq:potential}.

For (2), note that
\[
 b_{n,m}(r)=\sqrt{\frac{\binom{n+m}{n}}{m!}}\,r^m .
\]
If $r>0$, put $t=r/(r+\sqrt n)\in(0,1)$. Then
\[
 \frac{\binom{n+m}{n}r^{2m}}{m!}
 =\left(\binom{n+m}{n}t^m\right)\frac{(r^2/t)^m}{m!}.
\]
Cauchy--Schwarz, the binomial series
$\sum_{m\geq0}\binom{n+m}{n}t^m=(1-t)^{-(n+1)}$ and the exponential series
give
\begin{align*}
 \sum_{m\geq0}b_{n,m}(r)
 &\leq\left(\sum_{m\geq0}\binom{n+m}{n}t^m\right)^{1/2}
      \left(\sum_{m\geq0}\frac{(r^2/t)^m}{m!}\right)^{1/2}\\
 &=(1-t)^{-(n+1)/2}\exp\left(\frac{r^2}{2t}\right).
\end{align*}
Put $u=r/\sqrt n$. Since $\log(1+u)\leq u$ for $u\geq0$ and
$\sqrt n\geq1$,
\begin{align*}
 \frac{n+1}{2}\log(1+u)+\frac{r(r+\sqrt n)}{2}
 &\leq\frac{(n+1)r}{2\sqrt n}+\frac{r^2+r\sqrt n}{2}\\
 &=r\sqrt n+\frac{r}{2\sqrt n}+\frac{r^2}{2}\\
 &\leq r\sqrt n+\frac{r+r^2}{2}.
\end{align*}
Adding $\tfrac12\log(n!)$ gives (2) for $r>0$; the case $r=0$ is immediate
because $S_n(0)=\sqrt{n!}$. Finally $|\xi_{n+m}|\leq1$, so
\eqref{eq:derivative} and the triangle inequality give
$|F_n(z)|\leq S_n(|z|)$. Applying the bound just proved with $r=|z|$, and
using that $r\mapsto(r+r^2)/2$ is nondecreasing, gives \eqref{eq:upper}
whenever $|z|\leq R$.
\end{proof}

The second ingredient is anti-concentration. If $\xi$ is uniform on
$\overline\D$ then, for $a\neq0$, $w\in\C$ and $\varepsilon\geq0$,
\begin{equation}\label{eq:smallball}
 \Pp\{|a\xi+w|<\varepsilon\}\leq\left(\frac{\varepsilon}{|a|}\right)^2,
\end{equation}
because the set $\{\zeta:|a\zeta+w|<\varepsilon\}$ is the disk of centre
$-w/a$ and radius $\varepsilon/|a|$, whose normalised area is at most
$(\varepsilon/|a|)^2$.

\begin{lemma}[Pointwise logarithmic tail]\label{lem:pointwise}
Let $n\geq1$ and let $z\in\C$ satisfy $1\leq|z|\sqrt n$ and $|z|\leq\sqrt n$.
Then, for every $s\geq0$,
\begin{equation}\label{eq:pointwise-tail}
 \Pp\bigl\{p_n(z)-\log|F_n(z)|>s+L_n\bigr\}\leq e^{-2s}.
\end{equation}
\end{lemma}

\begin{proof}
Put $m=\lfloor|z|\sqrt n\rfloor$ and $k=n+m$, and split \eqref{eq:derivative}
by isolating the coefficient $\xi_k$:
\[
 F_n(\omega,z)=a\,\xi_k(\omega)+w(\omega),
\]
where
\[
 a=\frac{\sqrt{(n+m)!}}{m!}\,z^m
\]
is deterministic and
\[
 w(\omega)=\sum_{j\neq m}\xi_{n+j}(\omega)
 \frac{\sqrt{(n+j)!}}{j!}\,z^j
\]
is the remainder.
We first check that $w$ is finite and measurable. The series defining $w$
converges absolutely for every $\omega$, since $|\xi_{n+j}|\leq1$ and
$\sum_{j\geq0}a_{n,j}(|z|)=S_n(|z|)<\infty$ by
Lemma~\ref{lem:saddle}(2); so $w$ is finite everywhere, not merely almost
surely. It is also a measurable function of $\omega$, being a uniform
limit of continuous functions on $\Omega$. The independence needed below
comes from the fact that $w$ does not depend on the coordinate $\omega_k$:
it is measurable with respect to
$\mathcal G=\sigma(\xi_j:j\neq k)$, since the index $j=m$ is omitted from the
sum. Note also that $|z|\geq1/\sqrt n>0$, so $a\neq0$ and
$|a|=a_{n,m}(|z|)$.

Since $|z|=:r$ satisfies the hypotheses of Lemma~\ref{lem:saddle}(1), that
lemma gives $\log|a|\geq p_n(z)-L_n$. Hence, for $s\geq0$,
\[
 \bigl\{p_n(z)-\log|F_n(z)|>s+L_n\bigr\}
 =\bigl\{|F_n(z)|<e^{\,p_n(z)-L_n-s}\bigr\}
 \subseteq\bigl\{|a\xi_k+w|<|a|e^{-s}\bigr\}.
\]
Now write $\Omega=\overline\D\times\Omega'$, where the first factor carries
the coordinate $\omega_k$ and $\Omega'=\overline\D^{\,\N\setminus\{k\}}$
carries the remaining coordinates, and correspondingly
$\Pp=m_\D\otimes\Pp'$. Because $w$ is $\mathcal G$-measurable it is a
function of $\omega'\in\Omega'$ alone, so Tonelli's theorem gives
\[
 \Pp\bigl\{|a\xi_k+w|<|a|e^{-s}\bigr\}
 =\int_{\Omega'}m_\D\bigl\{\zeta\in\overline\D:
   |a\zeta+w(\omega')|<|a|e^{-s}\bigr\}\,d\Pp'(\omega').
\]
For each fixed $\omega'$ the inner measure is at most $(e^{-s})^2=e^{-2s}$ by
\eqref{eq:smallball}, with $\varepsilon=|a|e^{-s}$; the bound is uniform in
$\omega'$, so integrating over $\Omega'$ preserves it. This proves
\eqref{eq:pointwise-tail}.
\end{proof}

The pointwise estimate is used through the following circle-average bound.

\begin{lemma}[Circle-average upper tail]\label{lem:circle-tail}
Let $Y\geq0$ be jointly measurable on $\Omega\times[0,2\pi]$ and suppose that
\[
 \Pp\{Y(\cdot,\theta)>s\}\leq e^{-s/2},\qquad s\geq0,
\]
uniformly in $\theta$. Put
$\overline Y=(2\pi)^{-1}\int_0^{2\pi}Y(\cdot,\theta)\,d\theta$. Then
\[
 \Pp\{\overline Y\geq t\}\leq2e^{-t/8},\qquad t\geq8.
\]
\end{lemma}

\begin{proof}
The tail hypothesis and the layer-cake formula give
\[
 \E Y(\cdot,\theta)
 =\int_0^\infty\Pp\{Y(\cdot,\theta)>s\}\,ds\leq2
\]
for every $\theta$. Hence $\E\overline Y\leq2$ by Tonelli's theorem, and
$\overline Y$ is integrable. Put $E=\{\overline Y\geq t\}$ and $q=\Pp(E)$. Then
$t\ind_E\leq\overline Y\ind_E$, and Tonelli's theorem gives
\[
 tq\leq\E\bigl(\overline Y\ind_E\bigr)
 =\frac1{2\pi}\int_0^{2\pi}\E\bigl(Y(\cdot,\theta)\ind_E\bigr)\,d\theta.
\]
For every $\theta$, the layer-cake formula and
\[
 \Pp(\{Y(\cdot,\theta)>s\}\cap E)
 \leq\min\{q,e^{-s/2}\}
\]
yield
\begin{align*}
 \E\bigl(Y(\cdot,\theta)\ind_E\bigr)
 &\leq\int_0^\infty\min\{q,e^{-s/2}\}\,ds\\
 &=\int_0^{2\log(1/q)}q\,ds
   +\int_{2\log(1/q)}^\infty e^{-s/2}\,ds\\
 &=2q\log\frac eq,
\end{align*}
where the computation assumes $0<q\leq1$. If $q>0$ we may cancel $q$ and
obtain $t\leq2\log(e/q)=2+2\log(1/q)$, so $q\leq e^{1-t/2}$. For $t\geq8$ we
have $e^{1-t/2}=e^{1-3t/8}e^{-t/8}\leq e^{-2}e^{-t/8}\leq2e^{-t/8}$. The case
$q=0$ is immediate.
\end{proof}

\section{Zero-free disks and proof of the theorem}\label{sec:holes}

We first isolate the geometric gain that drives the estimate. For $c\in\C$
and $r>0$ set
\[
 \gamma(c,r)=\frac1{2\pi}\int_0^{2\pi}|c+re^{i\theta}|\,d\theta-|c|.
\]

\begin{lemma}\label{lem:gap}
For all $c\in\C$ and $r>0$ we have $\gamma(c,r)>0$.
\end{lemma}

\begin{proof}
Set
\[
 \varphi(\theta)=|c+re^{i\theta}|+|c-re^{i\theta}|-2|c|.
\]
The triangle inequality gives $\varphi\geq0$, and $\varphi$ is continuous. If
$c=0$, then $\varphi\equiv2r>0$. If $c\neq0$, choose $\theta_0$ with
$re^{i\theta_0}=irc/|c|$. The Pythagorean identity gives
\[
 \varphi(\theta_0)=2\sqrt{|c|^2+r^2}-2|c|>0.
\]
Thus continuity implies
$\int_0^{2\pi}\varphi(\theta)\,d\theta>0$. The substitution
$\theta\mapsto\theta+\pi$ gives
\[
 2\gamma(c,r)
 =\frac1{2\pi}\int_0^{2\pi}\varphi(\theta)\,d\theta>0,
\]
as required.
\end{proof}

Next we record the mean value step in the exact form needed, with the strict
radius margin made explicit.

\begin{lemma}[Mean value property on a zero-free disk]\label{lem:mvp}
Let $c\in\C$ and $\rho>0$, let $G$ be holomorphic and zero-free on
$B(c,\rho)$, and let $0<r<\rho$. Then
\[
 \frac1{2\pi}\int_0^{2\pi}\log|G(c+re^{i\theta})|\,d\theta=\log|G(c)|.
\]
\end{lemma}

\begin{proof}
The disk $B(c,\rho)$ is convex, and $G'/G$ is holomorphic on it because $G$ is
zero-free there. Every holomorphic function on a convex domain has a
holomorphic primitive \cite[Chapter~IV]{Conway}, so there is a holomorphic
$h$ on $B(c,\rho)$ with $h'=G'/G$. Then
$(Ge^{-h})'=e^{-h}(G'-Gh')=0$ on the connected set $B(c,\rho)$, so
$Ge^{-h}\equiv\kappa$ for a constant $\kappa$, and $\kappa\neq0$ because $G$
is zero-free. Choose $\lambda\in\C$ with $e^{\lambda}=\kappa$ and set
$g=h+\lambda$; then $e^{g}=G$ on $B(c,\rho)$, and hence
$\Real g=\log|G|$ there.

The strict inequality $r<\rho$ gives $\overline{B(c,r)}\subset B(c,\rho)$, so
the circle $|z-c|=r$ and its interior lie in the domain of $g$ and Cauchy's
integral formula \cite[Chapter~IV]{Conway} applies:
\[
 g(c)=\frac1{2\pi i}\int_{|z-c|=r}\frac{g(z)}{z-c}\,dz
 =\frac1{2\pi}\int_0^{2\pi}g(c+re^{i\theta})\,d\theta,
\]
the second equality by the parametrisation $z=c+re^{i\theta}$,
$dz=ire^{i\theta}\,d\theta$. Taking real parts gives the assertion.
\end{proof}

Now fix $c\in\Q+i\Q$ with $c\neq0$ and a rational $R$ with $0<R<|c|$, and put
\begin{equation}\label{eq:margin}
 r=\frac R2<R,
\end{equation}
so that $\overline{B(c,r)}\subset B(c,R)$ with a strict radius margin. On the
circle $|z-c|=r$ we have
\[
 0<\delta:=|c|-r\leq|z|\leq|c|+r=:M,
\]
the left inequality because $r<R<|c|$. Consequently there is an integer
$n_0=n_0(c,R)\geq1$ such that
\begin{equation}\label{eq:largen}
 1\leq\delta\sqrt n\quad\text{and}\quad M\leq\sqrt n
 \qquad\text{for all }n\geq n_0,
\end{equation}
and then the hypotheses of Lemma~\ref{lem:pointwise} hold at every point of
the circle $|z-c|=r$, uniformly in $\theta$.

To obtain a real-valued measurable deficit without changing it on the
zero-free event used below, define
\[
 \ell(w)=\begin{cases}\log|w|,&w\neq0,\\ 0,&w=0,\end{cases}
 \qquad
 D_{n,z}=\max\{p_n(z)-\ell(F_n(z))-L_n,\,0\},
\]
and, for $n\geq n_0$,
\[
 \overline D_n=\frac1{2\pi}\int_0^{2\pi}D_{n,c+re^{i\theta}}\,d\theta .
\]
The function $\ell$ is Borel measurable on $\C$, and $(\omega,z)\mapsto
F_n(\omega,z)$ is jointly continuous, so
$(\omega,\theta)\mapsto D_{n,c+re^{i\theta}}(\omega)$ is jointly measurable
on $\Omega\times[0,2\pi]$ and $\overline D_n$ is a measurable function of
$\omega$.

Because $\ell(w)\geq\log|w|$ for every $w\in\C$ in the
extended-real order, with equality unless $w=0$, we have, for $s\geq0$,
\[
 \{D_{n,z}>s\}\subseteq
 \{p_n(z)-\log|F_n(z)|>s+L_n\}.
\]
Combining this with Lemma~\ref{lem:pointwise} and using $e^{-2s}\leq e^{-s/2}$
for $s\geq0$ gives the tail hypothesis of Lemma~\ref{lem:circle-tail}
explicitly: for $n\geq n_0$, every $z$ on the circle $|z-c|=r$ and every
$s\geq0$,
\begin{equation}\label{eq:bridge}
 \Pp\{D_{n,z}>s\}
 \leq\Pp\{p_n(z)-\log|F_n(z)|>s+L_n\}
 \leq e^{-2s}
 \leq e^{-s/2}.
\end{equation}
So Lemma~\ref{lem:circle-tail} applies to $Y(\omega,\theta)
=D_{n,c+re^{i\theta}}(\omega)$ and to $\overline Y=\overline D_n$.

Let
\[
 H_n(c,R)=\{\omega\in\Omega:F_n(\omega,\cdot)\text{ has no zero in }B(c,R)\}.
\]

\begin{lemma}\label{lem:measurable}
$H_n(c,R)\in\mathcal F$.
\end{lemma}

\begin{proof}
We claim that
\begin{equation}\label{eq:meas-identity}
 H_n(c,R)=\bigcap_{\substack{\rho\in\Q\\ 0<\rho<R}}
 \Bigl\{\omega:\min_{|z-c|\leq\rho}|F_n(\omega,z)|>0\Bigr\},
\end{equation}
the minimum being attained because $z\mapsto|F_n(\omega,z)|$ is continuous
and $\overline{B(c,\rho)}$ is compact. If $\omega\in H_n(c,R)$ and
$0<\rho<R$, then $\overline{B(c,\rho)}\subset B(c,R)$, so $|F_n(\omega,
\cdot)|$ is continuous and strictly positive on a compact set and its minimum
there is positive. Conversely, if the right-hand side contains $\omega$ and
$z\in B(c,R)$, choose a rational $\rho$ with $|z-c|\leq\rho<R$; then
$|F_n(\omega,z)|>0$. This proves \eqref{eq:meas-identity}, and the
intersection is countable.

It remains to check that, for fixed $\rho\in(0,R)$, the function
\[
 M_\rho(\omega)=\min_{|z-c|\leq\rho}|F_n(\omega,z)|
\]
is $\mathcal F$-measurable. Let $Q_\rho=(\Q+i\Q)\cap B(c,\rho)$, a countable
set that is dense in $\overline{B(c,\rho)}$: boundary points are approached
from inside the disk and then by points of $\Q+i\Q$. For fixed $\omega$ the function
$z\mapsto|F_n(\omega,z)|$ is continuous on the compact set
$\overline{B(c,\rho)}$, and a continuous function on a compact set has the
same infimum over the set as over any dense subset of it; hence
\[
 M_\rho(\omega)=\inf_{z\in Q_\rho}|F_n(\omega,z)| .
\]
For each fixed $z$ the map $\omega\mapsto|F_n(\omega,z)|$ is measurable, as
noted after \eqref{eq:derivative}, so $M_\rho$ is a countable infimum of
measurable functions and is therefore measurable. Consequently
$\{M_\rho>0\}\in\mathcal F$ for every rational $\rho\in(0,R)$, and
\eqref{eq:meas-identity} exhibits $H_n(c,R)$ as a countable intersection of
sets in $\mathcal F$.
\end{proof}

We can now estimate $\Pp(H_n(c,R))$. Fix $n\geq n_0$ and work on the event
$H_n(c,R)$, on which $F_n$ is holomorphic and zero-free on $B(c,R)$ by the
definition of that event. Since $r=R/2<R$ by \eqref{eq:margin},
Lemma~\ref{lem:mvp} applies with $G=F_n$ and $\rho=R$ and gives
\begin{equation}\label{eq:jensen}
 \frac1{2\pi}\int_0^{2\pi}\log|F_n(c+re^{i\theta})|\,d\theta
 =\log|F_n(c)| .
\end{equation}
Lemma~\ref{lem:saddle}(2), applied with $R$ replaced by $|c|$ and evaluated
at $z=c$, gives
\begin{equation}\label{eq:centre}
 \log|F_n(c)|\leq p_n(c)+C(c),\qquad C(c)=\frac{|c|+|c|^2}{2}.
\end{equation}
On $H_n(c,R)$ we have $F_n(z)\neq0$ for $z\in\overline{B(c,r)}$, so
$\ell(F_n(z))=\log|F_n(z)|$ there and therefore
$D_{n,z}\geq p_n(z)-\log|F_n(z)|-L_n$ on the circle $|z-c|=r$. Averaging over
that circle and using
\[
 \frac1{2\pi}\int_0^{2\pi}p_n(c+re^{i\theta})\,d\theta
 =\tfrac12\log(n!)+\sqrt n\cdot\frac1{2\pi}\int_0^{2\pi}|c+re^{i\theta}|\,d\theta
 =p_n(c)+\gamma(c,r)\sqrt n,
\]
together with \eqref{eq:jensen} and \eqref{eq:centre}, we obtain on
$H_n(c,R)$ that
\[
 \overline D_n
 \geq p_n(c)+\gamma(c,r)\sqrt n-\log|F_n(c)|-L_n
 \geq\gamma(c,r)\sqrt n-C(c)-L_n=:T_n,
\]
that is,
\begin{equation}\label{eq:inclusion}
 H_n(c,R)\subseteq\{\overline D_n\geq T_n\},
 \qquad
 T_n=\gamma(c,r)\sqrt n-C(c)-2-\log(n+1).
\end{equation}
By Lemma~\ref{lem:gap} we have $\gamma(c,r)>0$, and
$C(c)+2+\log(n+1)=o(\sqrt n)$, so there is $n_1=n_1(c,R)\geq n_0$ with
\begin{equation}\label{eq:Tn}
 T_n\geq\max\Bigl\{\tfrac12\gamma(c,r)\sqrt n,\;8\Bigr\},
 \qquad n\geq n_1 .
\end{equation}
For $n\geq n_1$, Lemma~\ref{lem:circle-tail} together with
\eqref{eq:bridge}, \eqref{eq:inclusion} and \eqref{eq:Tn} gives
\begin{equation}\label{eq:hole}
 \Pp\bigl(H_n(c,R)\bigr)\leq\Pp\{\overline D_n\geq T_n\}
 \leq2e^{-T_n/8}
 \leq2\exp\left(-\frac{\gamma(c,r)}{16}\sqrt n\right).
\end{equation}
The right-hand side of \eqref{eq:hole} is summable, since
$\exp(-\gamma(c,r)\sqrt n/16)\leq n^{-2}$ for all large $n$, and the finitely
many terms with $n<n_1$ are each at most $1$. Hence
\begin{equation}\label{eq:summable-holes}
 \sum_{n=0}^{\infty}\Pp\bigl(H_n(c,R)\bigr)<\infty .
\end{equation}
All the events involved are measurable by Lemma~\ref{lem:measurable}, so the
first Borel--Cantelli lemma \cite[Section~2.3]{Durrett} applies and shows
that
\begin{equation}\label{eq:bc}
 \Pp\Bigl(\limsup_{n\to\infty}H_n(c,R)\Bigr)=0;
\end{equation}
that is, almost surely only finitely many $F_n$ are zero-free in $B(c,R)$.

\begin{proof}[Proof of Theorem~\ref{thm:main}]
Let $\mathcal B$ be the family of disks $B(c,R)$ with $c\in\Q+i\Q$,
$c\neq0$, $R\in\Q_{>0}$ and $R<|c|$; it is countable. We check that every
nonempty open set $U\subset\C$ contains a member of $\mathcal B$. Since $U$
is open and nonempty it contains a point $z_0\neq0$, and there is
$\varepsilon>0$ with $B(z_0,\varepsilon)\subset U$ and
$\varepsilon<|z_0|/2$. Choose $c\in\Q+i\Q$ with $|c-z_0|<\varepsilon/4$;
then $|c|\geq|z_0|-\varepsilon/4>0$, so $c\neq0$. Choose a rational $R$ with
$0<R<\min\{\varepsilon/2,|c|\}$. Then $R<|c|$, and every $z$ with $|z-c|<R$
satisfies $|z-z_0|<\varepsilon/4+\varepsilon/2<\varepsilon$, so
$B(c,R)\subset B(z_0,\varepsilon)\subset U$ and $B(c,R)\in\mathcal B$.

By \eqref{eq:bc}, for each $B=B(c,R)\in\mathcal B$ there is a
$\Pp$-null set $\mathcal N_B$ outside of which only finitely many $F_n$ are
zero-free in $B$. The union $\mathcal N=\bigcup_{B\in\mathcal B}\mathcal N_B$
is a countable union of null sets and is therefore null. With $\Omega_0$ as
in \eqref{eq:nonzero} we have $\Pp(\Omega_0\setminus\mathcal N)=1$; in
particular this event is nonempty. Fix $\omega\in\Omega_0\setminus\mathcal N$
and let $f$ be the entire function \eqref{eq:fock} determined by it.

Then $f$ is transcendental, because $\omega\in\Omega_0$ makes every Taylor
coefficient of $f$ nonzero, and $f$ satisfies \eqref{eq:growth}, which holds
for every sample. Let $U\subset\C$ be nonempty and open and choose
$B\in\mathcal B$ with $B\subset U$. Since $\omega\notin\mathcal N_B$, only
finitely many $F_n=f^{(n)}$ are zero-free in $B$. Let $N(U)$ be one more than
the largest such index, or let $N(U)=0$ if there is no such index. Then every
$n\geq N(U)$ gives a zero of $f^{(n)}$ in $B$, hence in $U$. This is the first
assertion of Theorem~\ref{thm:main}, and the last assertion follows from it by
Lemma~\ref{lem:cofinite}.
\end{proof}

\section{A consequence for a theorem of Boas and Reddy}
\label{sec:boas-reddy}

Under their standing assumption that $f$ is transcendental, Boas and Reddy
assert that if $f$ is of order at most two and finite type, then $f$ admits
arbitrarily large disks on which infinitely many successive derivatives are
zero-free \cite[Theorem~1]{BoasReddy1973}. They clarify that
one fixed disk serves infinitely many derivatives \cite[p.~65]{BoasReddy1973}.
Their expanded article also stated a distinct Theorem~2
\cite{BoasReddyJMAA1973}.
Their published erratum addresses Theorem~2 \cite{BoasReddy1975}.
It withdraws the supporting construction and says that conclusion remains
open.

\begin{corollary}
The function in Theorem~\ref{thm:main} is a counterexample to the assertion of
\cite[Theorem~1]{BoasReddy1973} as printed.
\end{corollary}

\begin{proof}
Let $f$ be the function supplied by Theorem~\ref{thm:main}. The estimate
\eqref{eq:growth} places $f$ in the growth class assumed by Boas and Reddy.
For every fixed nonempty disk $D$, Theorem~\ref{thm:main} supplies an integer
$N(D)$ such that $f^{(n)}$ has a zero in $D$ whenever $n\geq N(D)$.
Consequently only the finitely many derivatives with $n<N(D)$ can be
zero-free in $D$. Thus no disk, of any radius, is zero-free for infinitely
many successive derivatives of $f$, contrary to the printed conclusion.
\end{proof}

\section*{Declarations}
\noindent\textit{Use of AI.} During preparation, OpenAI's GPT-5.6 Sol assisted
with the small-ball estimate and editing. Anthropic's Claude Fable and Opus and
OpenAI Codex assisted with manuscript revision and argument auditing. The
author checked all AI-assisted output and takes sole responsibility for the
mathematical content and exposition of this article.

\noindent\textit{Competing interests.} The author declares no competing
interests. \textit{Funding.} The author received no specific funding for this
work. \textit{Data availability.} No data were generated or analysed in this
theoretical study.

\begin{thebibliography}{99}

\bibitem{Almeida2026}
A.~Almeida,
`A probabilistic construction for the transcendental form of Erd\H{o}s
Problem 906', unpublished manuscript (April 2026),
\href{https://github.com/Drill23/erdos-problem-906-proposed-solution/commit/5d77fdc959dc6df150d3c680b338ed7b99653673}
{GitHub repository, commit \texttt{5d77fdc}}, accessed 16 August 2026.

\bibitem{AlmeidaForum2026}
A.~Almeida,
`Proposed solution to Erd\H{o}s Problem 906', forum post, 25 April 2026,
\url{https://www.erdosproblems.com/forum/thread/906#post-5875},
accessed 16 August 2026.

\bibitem{BarthSchneider1972}
K.~F. Barth and W.~J. Schneider,
`On a problem of Erd\H{o}s concerning the zeros of the derivatives of an
entire function',
\emph{Proc. Amer. Math. Soc.} \textbf{32} (1972), 229--232.

\bibitem{Bloom906}
T.~F. Bloom,
`Erd\H{o}s Problem \#906',
\href{https://www.erdosproblems.com/906}{erdosproblems.com/906},
accessed 16 August 2026.

\bibitem{BoasReddy1973}
R.~P. Boas, Jr. and A.~R. Reddy,
`Zeros of successive derivatives of entire functions',
\emph{Bull. Amer. Math. Soc.} \textbf{79} (1973), 64--65,
\url{https://doi.org/10.1090/S0002-9904-1973-13093-9}.

\bibitem{BoasReddyJMAA1973}
R.~P. Boas, Jr. and A.~R. Reddy,
`Zeros of the successive derivatives of entire functions',
\emph{J. Math. Anal. Appl.} \textbf{42} (1973), 466--473,
\url{https://doi.org/10.1016/0022-247X(73)90153-4}.

\bibitem{BoasReddy1975}
R.~P. Boas, Jr. and A.~R. Reddy,
`Erratum',
\emph{J. Math. Anal. Appl.} \textbf{49}(2) (1975), 527,
\url{https://doi.org/10.1016/0022-247X(75)90195-X}.

\bibitem{ChojeckiForum2026}
P.~Chojecki,
`Posts on Erd\H{o}s Problem 906', 25 April and 1 June 2026,
\url{https://www.erdosproblems.com/forum/thread/906#post-5886} and
\url{https://www.erdosproblems.com/forum/thread/906#post-6776},
accessed 16 August 2026.

\bibitem{Conway}
J.~B. Conway,
\emph{Functions of One Complex Variable I}, 2nd edn,
(Springer, New York, 1978).

\bibitem{Durrett}
R.~Durrett,
\emph{Probability: Theory and Examples}, 5th edn,
(Cambridge University Press, 2019).

\bibitem{Erdos1982}
P.~Erd\H{o}s,
`Some of my favourite problems which recently have been solved',
in: \emph{Proc. Internat. Math. Conf., Singapore 1981},
North-Holland Math. Stud. 74 (North-Holland, 1982), 59--79,
\url{https://doi.org/10.1016/S0304-0208(08)70415-8}.

\bibitem{Kahane}
J.-P.~Kahane,
\emph{Some Random Series of Functions}, 2nd edn,
(Cambridge University Press, 1985).

\bibitem{MacLane1952}
G.~R. MacLane,
`Sequences of derivatives and normal families',
\emph{J. Anal. Math.} \textbf{2} (1952), 72--87.

\end{thebibliography}

\end{document}
